\section{Discussion}
From a translational perspective, this public study reinforces that
structural variation should not be treated as a secondary genomic layer
in pediatric oncology. Even when overall rearrangement burden is lower
than that in adult solid tumors, the disrupted loci can still
concentrate in biologically decisive drivers and therefore remain highly
relevant for disease classification and downstream clinical testing
\supercite{Greenhalgh_2026}.

The RAG-associated hotspot and signature findings are also important
conceptually because they connect genome architecture to
lineage-specific mutational processes. That link helps explain why
lymphoid malignancies may display recurrent rearrangement patterns that
differ from those seen in non-hematological tumors and provides a
clearer mechanistic narrative than burden summaries alone \supercite{Greenhalgh_2026}.

This repository example intentionally stops at the level of a concise
public summary. Users who need full methodological detail, extended
figure interpretation, or exact supplemental evidence should consult the
published article directly rather than rely on this demonstration text.
The value of the current project lies in showing how an openly shareable
manuscript skeleton can be rendered into publication-style PDF and DOCX
artifacts with a single build command.
